\documentclass[11pt]{article}

\usepackage[T1]{fontenc}
\usepackage[utf8]{inputenc}
\usepackage[margin=0.75in]{geometry}
\usepackage{amsmath,amssymb}
\usepackage{booktabs,longtable,array,tabularx}
\usepackage{enumitem}
\usepackage{xcolor}
\usepackage{hyperref}
\usepackage[nameinlink,noabbrev]{cleveref}
\usepackage[backend=biber,style=numeric,sorting=none,maxbibnames=10]{biblatex}
\addbibresource{refs.bib}
\usepackage[most]{tcolorbox}
\usepackage{microtype}

\setcounter{tocdepth}{2}
\setcounter{secnumdepth}{2}
\hypersetup{colorlinks=true,linkcolor=blue!60!black,citecolor=blue!60!black,urlcolor=blue!60!black}
\setlist[itemize]{topsep=2pt,itemsep=1pt,parsep=1pt}
\setlist[enumerate]{topsep=2pt,itemsep=1pt,parsep=1pt}
\renewcommand{\arraystretch}{1.15}

\newcolumntype{Y}{>{\raggedright\arraybackslash}X}
\newcolumntype{L}[1]{>{\raggedright\arraybackslash}p{#1}}
\newtcolorbox{decisionbox}{colback=green!3!white,colframe=green!40!black,title=Recommended Decision,breakable}
\newtcolorbox{warningbox}{colback=orange!5!white,colframe=orange!60!black,title=Important Constraint,breakable}
\newtcolorbox{scopebox}{colback=blue!3!white,colframe=blue!50!black,title=Scope and Currency,breakable}

\title{Engineering Hardware and Local AI Trade Study\\
\large NavKit Development, Local Inference, Simulation, and Future Perception Workloads}
\author{NavKit}
\date{Planning baseline: 9 October 2026}

\begin{document}
\hypersetup{pageanchor=false}
\maketitle
\hypersetup{pageanchor=true}
\tableofcontents
\clearpage

\section{Executive Guidance}

\begin{scopebox}
This document is written for a human purchasing decision, not as an autonomous
agent handoff. Dollar figures are approximate United States planning bands as
of the date on the title page. Manufacturer list prices, used-market estimates,
and complete-system budgets are deliberately distinguished. Obtain current
quotes before purchasing.
\end{scopebox}

The recommended architecture is two independent machines:

\begin{enumerate}
  \item a durable portable workstation for editing, compilation, debugging,
        visualization, and field work; and
  \item an expandable Linux server for local LLM inference first, followed by
        Monte Carlo campaigns, continuous integration, data storage, vision
        processing, and selected factor-graph workloads.
\end{enumerate}

\begin{decisionbox}
Purchase the laptop first. Target a used ThinkPad P16 Gen 1 or Gen 2 with an HX
processor, at least 64~GB of RAM, and two NVMe slots. Do not pay a large premium
for its GPU.

For the server, either buy a deliberately inexpensive 16~GB GPU as a learning
device or wait until a 24--32~GB GPU fits the budget. Build the surrounding
platform for the expected final server: sufficient slot spacing, power,
cooling, memory, storage, and networking. Do not buy Threadripper until measured
CPU scaling or a concrete PCIe/memory requirement justifies it.
\end{decisionbox}

The near-term sequence is:

\begin{enumerate}
  \item acquire and validate the laptop immediately;
  \item capture baseline NavKit build, simulation, packaging, and plotting
        measurements;
  \item begin local inference in early 2027 with a 16~GB entry GPU if necessary;
  \item upgrade to or add a 24--32~GB GPU when model capacity or vision work
        demonstrates the need; and
  \item consider 48~GB professional GPUs, multiple GPUs, or Threadripper only
        after a measured workload crosses the relevant boundary.
\end{enumerate}

\section{Workload Model}

\subsection{What uses which hardware}

\begin{table}[htbp]
\centering
\small
\begin{tabularx}{\textwidth}{L{0.25\textwidth} L{0.23\textwidth} Y}
\toprule
Workload & Dominant resource & Hardware implication \\
\midrule
Interactive C++ development & Fast CPU cores, RAM, NVMe & Laptop responsiveness matters more than a laptop GPU. \\
Clean parallel build & CPU cores, RAM, cooling & Sustained power and 64+~GB avoid contention. \\
Single SWIL run & CPU latency & Favor modern per-core performance. \\
Monte Carlo campaigns & CPU processes, RAM, storage & More cores help independent runs; benchmark scaling before buying HEDT. \\
HDF5 packaging and analysis & NVMe, RAM, CPU & Fast local scratch storage matters; a GPU currently adds little. \\
Plotly/Matplotlib rendering & CPU, RAM & Cached HDF5 already changes performance by an order of magnitude. \\
Local LLM inference & GPU VRAM, memory bandwidth & Capacity determines which model and context fit; compute determines response rate. \\
Vision navigation & GPU, NVMe, CPU & Dataset decoding, learned front ends, and bundle adjustment create mixed demand. \\
Factor graphs & CPU first; GPU for large supported graphs & GPU benefit depends on graph size, solver, and retained GPU residency. \\
HWIL orchestration & Deterministic CPU/network timing & Isolate from LLM inference and nondeterministic batch jobs. \\
Artifact/data storage & Capacity, integrity, backup & Treat scratch, evidence, and backup as separate roles. \\
\bottomrule
\end{tabularx}
\caption{Resource map for the intended engineering workloads.}
\end{table}

\subsection{Evidence from NavKit}

NavKit's current implementation is CPU-oriented C++23 and Eigen estimation and
simulation with Python/HDF5 analysis. Existing measurements show that:

\begin{itemize}
  \item a representative 1,000-run plot took 213.198~s from CSV, 10.275~s from
        cached HDF5 with Matplotlib, and 16.868~s from HDF5 with Plotly;
  \item the full 1,000-run analysis bundle was 2.32~GB and initially required
        1,319.129~s to package;
  \item four analysis workers improved complete 500- and 1,000-run workloads by
        only 11.5\% and 16.8\% over serial execution; and
  \item later cache work reduced 1,000-run cold preparation from 209.189~s to
        75.145~s, after which HDF5 input became the primary remaining cost.
\end{itemize}

These measurements do not justify purchasing a GPU for present NavKit
analysis. They favor strong single-core performance, enough CPU cores for
independent simulations, 64--128~GB of RAM, and fast NVMe storage. GPU value
arrives through local inference and future perception/smoothing work.

\section{Portable Workstation}

\subsection{Option matrix}

\begin{longtable}{L{0.16\textwidth} L{0.14\textwidth} L{0.16\textwidth} L{0.17\textwidth} L{0.17\textwidth}}
\toprule
Option & Planning band & Advantages & Limitations & Guidance \\
\midrule
\endfirsthead
\toprule
Option & Planning band & Advantages & Limitations & Guidance \\
\midrule
\endhead
Used P16 Gen 1 & \$900--\$1,350 & Four SODIMMs, up to 128~GB, two NVMe slots, HX CPU, serviceable & Heavy; older efficiency and battery life & Best value if condition, battery, thermals, and display are verified. \\
Used P16 Gen 2 & \$1,300--\$2,000 & Newer HX CPUs; up to 192~GB non-ECC or 128~GB ECC; two NVMe slots & Higher cost; still a transportable workstation & Prefer when the premium over Gen 1 is modest. \\
New mobile workstation & \$2,500+ & Warranty, modern efficiency, known history & Rapid depreciation; expensive GPU upgrades; soldered-memory risk on some models & Buy only for warranty, battery life, or employer-style support requirements. \\
Thin performance laptop & \$1,500--\$2,500 & Better portability and battery life & Reduced sustained cooling, RAM, repairability, and storage expansion & Poorer match for this development workflow. \\
\bottomrule
\caption{Portable workstation alternatives. Price bands are planning estimates, not quotes.}
\end{longtable}

Lenovo specifies four memory slots and up to 128~GB for the P16 Gen 1
\cite{lenovo_p16_gen1}. The P16 Gen 2 supports two NVMe drives and up to 192~GB
non-ECC or 128~GB ECC depending on configuration \cite{lenovo_p16_gen2}.

\subsection{Recommended laptop configuration}

\begin{table}[htbp]
\centering
\small
\begin{tabularx}{\textwidth}{L{0.21\textwidth} L{0.27\textwidth} Y}
\toprule
Subsystem & Target & Rationale \\
\midrule
CPU & Intel HX, favoring condition and price over the final SKU step & Good interactive and parallel compile performance. \\
RAM & 64~GB minimum; 128~GB preferred if the incremental cost is reasonable & Supports IDE, WSL/VMs, builds, analysis, and large HDF5 products concurrently. \\
Storage & 2~TB primary NVMe; retain the second slot for expansion & Avoids immediate storage pressure while keeping datasets replaceable. \\
GPU & Integrated or modest discrete GPU & The server owns serious inference and future vision compute. \\
Display & Verify resolution, brightness, scaling, and panel defects in person & Affects every development hour. \\
Condition & Battery health, hinges, keyboard, ports, fans, charger, BIOS access & More important than a small CPU-bin difference in a used workstation. \\
Warranty & Remaining Lenovo coverage or priced replacement risk & A used workstation is still a field-development dependency. \\
\bottomrule
\end{tabularx}
\caption{Recommended portable workstation priorities.}
\end{table}

\section{Server Platform}

\subsection{CPU and platform matrix}

Budgets in this section exclude the GPU but include a plausible motherboard,
CPU, cooling, memory, chassis, PSU, and initial storage.

\begin{longtable}{L{0.16\textwidth} L{0.14\textwidth} L{0.16\textwidth} L{0.20\textwidth} L{0.16\textwidth}}
\toprule
Platform & Base-system band & Strength & Constraint & Best use \\
\midrule
\endfirsthead
\toprule
Platform & Base-system band & Strength & Constraint & Best use \\
\midrule
\endhead
Used AM4 / Ryzen 5950X & \$750--\$1,250 & Cheap 16-core throughput; mature DDR4; strong entry campaign runner & End-of-line platform; dual-channel memory; limited PCIe expansion & Intentional budget server with one primary GPU. \\
Modern AM5 / Ryzen 9950X class & \$1,400--\$2,300 & 16 modern cores, strong single-thread performance, DDR5, long-lived mainstream platform & Consumer PCIe lanes and board layouts constrain serious dual-GPU builds & Recommended balanced single-GPU server. \\
Used workstation platform & \$1,800--\$3,500 & More PCIe lanes, ECC capacity, remote-management options & Older per-core performance, power, noise, and uncertain history & Good when expansion is worth more than peak interactive speed. \\
Threadripper 9960X class & \$3,500--\$6,000 & 24 cores, quad-channel RDIMM, ECC, 88 usable PCIe lanes & 350~W CPU; costly board, memory, and cooling & Multiple accelerators, high I/O, or proven CPU scaling. \\
\bottomrule
\caption{Server CPU/platform alternatives.}
\end{longtable}

AMD specifies 16 cores, 32 threads, and a 170~W TDP for the Ryzen 9 9950X
class \cite{amd_9950x}. The Threadripper 9960X provides 24 cores, quad-channel
RDIMM ECC memory, and 88 usable PCIe lanes at a 350~W default TDP
\cite{amd_9960x}. Those extra lanes and memory channels are valuable only when
the installed devices or measured workload use them.

\begin{decisionbox}
For a server expected to remain primarily single-GPU, use a modern 16-core AM5
platform with 128~GB of RAM. If two large GPUs are a credible near-term plan,
select the exact motherboard and chassis from a physical lane/slot drawing---or
move directly to a workstation platform. Do not assume that two nominal x16
slots are electrically suitable, correctly spaced, or simultaneously usable
with the required NVMe and networking devices.
\end{decisionbox}

\subsection{Memory matrix}

\begin{table}[htbp]
\centering
\small
\begin{tabularx}{\textwidth}{L{0.14\textwidth} L{0.20\textwidth} Y L{0.20\textwidth}}
\toprule
Capacity & Suitable for & Trade & Recommendation \\
\midrule
64~GB & Entry inference server, ordinary builds, modest campaigns & Can become constrained by VMs, HDF5, datasets, and concurrent services & Laptop minimum; server only under a hard budget. \\
128~GB & NavKit campaigns, CI, analysis, local inference, containers & Best mainstream balance & Recommended server baseline. \\
192--256~GB & Large CPU factor graphs, many services, RAM caching, large datasets & Higher DIMM/platform cost; possible memory-speed reduction & Buy only with a defined consumer. \\
512~GB+ & Large in-memory graphs or shared laboratory infrastructure & Workstation/server platform required & Deferred until measurements require it. \\
\bottomrule
\end{tabularx}
\caption{System-memory capacity trade.}
\end{table}

\section{GPU and Local Inference}

\subsection{GPU option matrix}

VRAM capacity is the first-order constraint for local inference. Compute and
bandwidth affect response speed only after the desired model, context, and cache
fit. The prices below distinguish manufacturer launch references from used or
professional-market planning bands; actual 2026 street prices may be much
higher.

\begin{longtable}{L{0.15\textwidth} L{0.09\textwidth} L{0.13\textwidth} L{0.18\textwidth} L{0.25\textwidth}}
\toprule
GPU class & VRAM & Price reference & Strength & Limitation / decision \\
\midrule
\endfirsthead
\toprule
GPU class & VRAM & Price reference & Strength & Limitation / decision \\
\midrule
\endhead
RTX 5060 Ti 16~GB & 16~GB & \$429 launch MSRP & Lowest sensible new NVIDIA entry; efficient; current CUDA stack & Limited model/context headroom. Buy only as an intentional starter. \\
Used RTX 3090 & 24~GB & Roughly \$600--\$1,000 used planning band & Excellent capacity per used-market dollar; practical 30--32B quantized models & 350~W, heat, age, no warranty, and uncertain prior use. Inspect carefully. \\
RTX 5080 & 16~GB & \$999 launch MSRP & Much more compute and bandwidth than entry cards & Poor LLM capacity value: still 16~GB. Favor only when vision throughput dominates. \\
RTX 5090 & 32~GB & \$1,999 launch MSRP; street quote required & Strong single-card inference/vision tier; materially more context and model capacity & Very high board power, physical size, price, and availability premium. \\
RTX PRO 5000 Blackwell & 48 or 72~GB & Professional quote & ECC, dual-slot form, 300~W, large contiguous memory, professional support & Several times consumer-GPU cost; defer until 32~GB is proven insufficient. \\
Two consumer GPUs & Varies & Sum of cards plus platform & Excellent for separate jobs; some engines can shard a model & Not transparent unified VRAM; PCIe communication and software support matter. \\
\bottomrule
\caption{GPU alternatives for local inference and future engineering compute.}
\end{longtable}

NVIDIA launched the 16~GB RTX 5060 Ti at \$429 \cite{nvidia_5060ti}. The RTX
5080 has 16~GB and a \$999 launch price \cite{nvidia_5080}; this makes it a
compute upgrade, not a model-capacity upgrade. NVIDIA specifies 32~GB for the
RTX 5090 \cite{nvidia_5090}. The professional RTX PRO 5000 Blackwell provides
48 or 72~GB of ECC memory in a dual-slot, 300~W design
\cite{nvidia_rtx_pro_5000}.

\begin{warningbox}
Two 16~GB GPUs are not universally equivalent to one 32~GB GPU. They provide
32~GB aggregate capacity only when the inference engine explicitly shards the
model and its runtime data. They are excellent for two independent model
servers, two agents, or concurrent vision and analysis jobs. Tensor-parallel
inference across consumer cards pays PCIe communication costs, and GeForce
cards do not provide a general transparent shared-memory pool.
\end{warningbox}

\subsection{Approximate model-capacity matrix}

For a model with $P$ parameters stored at $b$ bits per parameter, weights alone
require approximately

\begin{equation}
  M_{\mathrm{weights}} \approx \frac{P b}{8}.
\end{equation}

Inference additionally needs quantization metadata, runtime workspaces, and a
key/value cache that grows with context length, batch size, layer count, and
hidden dimension. Consequently, weight-only arithmetic is a lower bound.

\begin{table}[htbp]
\centering
\small
\begin{tabularx}{\textwidth}{L{0.13\textwidth} L{0.25\textwidth} Y L{0.23\textwidth}}
\toprule
VRAM & Comfortable target & Possible with compromise & Engineering interpretation \\
\midrule
16~GB & 7--14B quantized models; modest VLMs & Some 20--32B models with aggressive quantization, short context, or CPU offload & Good learning tier; expect capacity decisions. \\
24~GB & 14B models with generous context; many 30--32B Q4 uses & Larger contexts may require tuning/offload & Strong used-market local-development tier. \\
32~GB & 30--32B quantized models with useful headroom & Selected larger low-bit models & Recommended durable single-consumer-GPU tier. \\
48~GB & Large 30--32B deployments; some 70B low-bit use & 70B context/batch still requires care & Professional capacity threshold. \\
72--96~GB & Large models, contexts, VLMs, or multiple services & Cost, not basic fit, becomes dominant & Laboratory infrastructure, not an entry purchase. \\
\bottomrule
\end{tabularx}
\caption{Conservative qualitative model-fit guidance. Exact fit depends on architecture and inference engine.}
\end{table}

\subsection{Recommended GPU path}

\begin{enumerate}
  \item If the early-2027 budget is constrained, buy a 16~GB RTX 5060 Ti only
        if its current street price remains proportionate to its \$429 launch
        reference. Use it to validate Linux, drivers, containers, serving,
        quantization, remote access, and actual model quality.
  \item If a reliable used RTX 3090 is available near the entry card's actual
        street price, its 24~GB capacity may be more useful for LLM work. Accept
        the energy, cooling, condition, and warranty risks explicitly.
  \item Prefer a future 32~GB single GPU over two 16~GB GPUs when the primary
        requirement is one larger model. Prefer two GPUs when independent jobs
        and service isolation are the primary requirements.
  \item Retain the starter GPU after an upgrade only if the platform can house
        and power it without compromising the primary GPU. It can then serve a
        small always-on model, vision preprocessing, or separate agent.
\end{enumerate}

\section{Storage, Network, Chassis, and Power}

\subsection{Storage tiers}

\begin{table}[htbp]
\centering
\small
\begin{tabularx}{\textwidth}{L{0.19\textwidth} L{0.22\textwidth} Y L{0.20\textwidth}}
\toprule
Role & Initial target & Contents & Policy \\
\midrule
OS/build/cache NVMe & 2~TB high-endurance TLC & Linux, source trees, containers, compiler cache, model runtime & Fast and replaceable; do not make it the only evidence store. \\
Scratch/dataset NVMe & 4~TB TLC & Monte Carlo output, HDF5, vision datasets, model weights & Optimize capacity and sustained write behavior. \\
Bulk artifact storage & 8--20~TB mirrored HDD or separate NAS & Retained campaigns, raw captures, releases, archived models & Capacity tier; not a substitute for backup. \\
Backup & Separate physical system or encrypted off-site target & Source-independent datasets, configs, evidence, documents & Test restoration; RAID alone is not backup. \\
\bottomrule
\end{tabularx}
\caption{Storage roles should be explicit rather than collapsed onto one disk.}
\end{table}

Avoid low-endurance QLC drives for the primary high-churn analysis scratch tier
unless their measured sustained-write behavior is acceptable. Record datasets
and qualification artifacts should be content-addressed or revision-labeled so
that a model or analysis result remains traceable to source, configuration,
seed, and tool version.

\subsection{Network, chassis, cooling, and UPS}

\begin{table}[htbp]
\centering
\small
\begin{tabularx}{\textwidth}{L{0.16\textwidth} L{0.23\textwidth} Y}
\toprule
Subsystem & Minimum / preferred & Guidance \\
\midrule
Network & 2.5~GbE / 10~GbE-ready & 2.5~GbE is adequate initially; 10~GbE becomes valuable for large vision and HDF5 datasets. \\
Chassis & Full tower or 4U+ rack chassis & Verify GPU length, height, width, cable bend radius, intake area, and the second-slot layout from drawings. \\
PSU & Quality ATX 3.x, sized with transient margin & A single 16~GB card may need 850--1,000~W; a 5090-class or dual-GPU plan commonly warrants 1,200--1,600~W after calculation. \\
Cooling & Front-to-back airflow with replaceable filters & Favor sustained compute and acoustics over decorative density. \\
UPS & USB/network-managed pure-sine unit & Size for controlled shutdown, not hours of full-load inference. \\
Remote recovery & Auto power-on, smart PDU/plug, or server BMC & Consumer platforms lack full out-of-band management; plan recovery before running unattended jobs. \\
\bottomrule
\end{tabularx}
\caption{Infrastructure choices that determine whether later GPU upgrades are practical.}
\end{table}

\section{Engineering Compute Beyond LLM Inference}

\subsection{Factor graphs}

Factor-graph optimization is a numerical workload, not an LLM workload. The
same GPU can serve both at different times, but the software stacks and
performance criteria are distinct.

GTSAM 4.3 introduced experimental CUDA acceleration for general nonlinear
optimization and structure from motion \cite{gtsam_43}. In the general path,
factor linearization, nonlinear error evaluation, and retraction remain on the
CPU while the sparse linear solve runs on the GPU. The project reports that
small problems can remain faster on the CPU; larger A100 benchmarks achieved
up to 3.64x for 2-D pose graphs, 4.70x for 3-D pose graphs, and 3.58x for stereo
SLAM. Its specialized GPU-resident structure-from-motion path reported roughly
8--10x acceleration on the tested BAL problems \cite{gtsam_cuda}.

The purchasing implication is to begin factor-graph work on the available CPU,
then benchmark representative graph sizes and solver choices. Do not purchase
a GPU or HEDT platform solely from the existence of a CUDA backend.

\subsection{Vision navigation}

Vision makes the strongest future case for a large GPU. Candidate uses include
learned features, optical flow, stereo/depth, image rectification pipelines,
visual place recognition, bundle adjustment, visual-inertial smoothing, map
construction, and multimodal inspection. It also increases storage and network
demand: model weights are small relative to retained image sequences and
derived datasets.

A 16~GB GPU is adequate for learning and prototype inference. A 24--32~GB GPU
is the more durable target for concurrent image batches, larger visual models,
and optimization. The exact upgrade should follow a representative camera
resolution, frame rate, model, and latency benchmark.

\subsection{SWIL, CI, services, and HWIL}

The server can grow into a useful private engineering platform:

\begin{itemize}
  \item scheduled NavKit builds, tests, and qualification campaigns;
  \item Monte Carlo execution and cached analysis generation;
  \item self-hosted development applications before cloud deployment;
  \item Git mirrors, package caches, model registry, and artifact storage;
  \item replay and mixed-HWIL orchestration; and
  \item later external-flight-computer stimulus and evidence collection.
\end{itemize}

Real-time or timing-sensitive HWIL must be isolated from LLM inference. Use
separate services and containers, explicit CPU affinity where necessary, resource
limits, a GPU job queue, and measured latency/jitter. A convenient shared server
must not silently turn model inference into a source of nondeterministic test
timing.

\section{Acquisition Paths}

\subsection{Path A: minimum-cost learning}

\begin{table}[htbp]
\centering
\small
\begin{tabularx}{\textwidth}{L{0.23\textwidth} L{0.20\textwidth} Y}
\toprule
Stage & Budgetary target & Outcome \\
\midrule
Laptop now & \$1,000--\$1,700 & P16 Gen 1/2, 64~GB, 2~TB; productive development immediately. \\
Server base & \$900--\$1,500 & Used AM4/5950X or carefully priced modern platform, 64--128~GB, 2~TB. \\
Starter GPU & Launch reference \$429; current quote required & 16~GB local inference and CUDA/vision experimentation. \\
Later upgrade & Workload-dependent & Replace the starter card or retain it for a separate small service. \\
\bottomrule
\end{tabularx}
\caption{Low-entry-cost path. It accepts a likely GPU replacement.}
\end{table}

\subsection{Path B: durable balanced server}

\begin{table}[htbp]
\centering
\small
\begin{tabularx}{\textwidth}{L{0.23\textwidth} L{0.20\textwidth} Y}
\toprule
Stage & Budgetary target & Outcome \\
\midrule
Laptop now & \$1,000--\$1,700 & Same portable workstation; server timing does not block development. \\
Server base & \$1,600--\$2,600 & Modern 16-core CPU, 128~GB, two NVMe drives, large chassis, upgrade-ready PSU. \\
Primary GPU & Market-dependent 24--32~GB tier & Useful 30--32B local models, more context, serious vision headroom. \\
Infrastructure & \$400--\$1,200 & UPS, networking, bulk storage, and backup according to existing equipment. \\
\bottomrule
\end{tabularx}
\caption{Preferred long-lived single-GPU engineering server.}
\end{table}

\subsection{Path C: workstation infrastructure}

Move to Threadripper or a workstation/server platform only after documenting at
least one of the following:

\begin{itemize}
  \item representative simulation throughput continues scaling beyond 16 CPU cores;
  \item required memory exceeds mainstream platform capacity or bandwidth;
  \item two or more full-bandwidth GPUs are required;
  \item storage, capture, and 10/25~Gb networking exhaust mainstream PCIe lanes; or
  \item several users or persistent services make manageability and ECC capacity
        more valuable than acquisition and operating cost.
\end{itemize}

\section{Benchmark and Purchase Gates}

\subsection{Laptop acceptance}

Before the return window closes:

\begin{enumerate}
  \item inspect battery health, thermals, fans, display, ports, keyboard, hinges,
        charger, BIOS access, and both storage slots;
  \item run memory and NVMe health tests;
  \item perform clean Debug and Release builds and capture elapsed time;
  \item run the complete test suite;
  \item run representative single, 100-, 500-, and 1,000-run NavKit workloads;
  \item record simulation, HDF5 packaging, and plotting time separately; and
  \item record peak RAM, CPU utilization, temperature, throttling, and disk throughput.
\end{enumerate}

\subsection{Starter-GPU acceptance}

Within the return window:

\begin{enumerate}
  \item validate the intended Linux distribution, NVIDIA driver, CUDA runtime,
        container runtime, and inference server;
  \item measure time to first token, sustained generation rate, VRAM use, power,
        temperature, and acoustics for selected 7--14B and larger quantized models;
  \item test the actual context length used for repository work rather than a
        short synthetic prompt;
  \item evaluate answer quality on a private NavKit task set; and
  \item run a small CUDA vision and GTSAM experiment to validate the broader stack.
\end{enumerate}

\subsection{Upgrade gates}

\begin{table}[htbp]
\centering
\small
\begin{tabularx}{\textwidth}{L{0.28\textwidth} Y}
\toprule
Observed limitation & Upgrade response \\
\midrule
Desired model or context does not fit 16~GB & Move to one 24--32~GB GPU; do not assume a second 16~GB card solves every case. \\
Model fits but is too slow & Upgrade GPU compute/bandwidth after profiling prompt processing and generation separately. \\
Independent services contend & Add a second GPU and assign services explicitly. \\
Monte Carlo CPU remains saturated and scales with more workers & Consider more CPU cores after controlling for storage and memory. \\
HDF5 or datasets dominate & Upgrade scratch NVMe, memory caching, or analysis format before CPU/GPU. \\
PCIe devices cannot coexist at required bandwidth & Move to Threadripper/workstation rather than stacking adapters onto a consumer board. \\
HWIL timing is perturbed & Isolate workloads or use a dedicated orchestration host; raw compute is not the remedy. \\
\bottomrule
\end{tabularx}
\caption{Measurements should trigger upgrades, not product enthusiasm.}
\end{table}

\section{Final Recommendation}

The preferred plan is:

\begin{enumerate}
  \item Buy a sound P16 Gen 1 or Gen 2 with 64~GB immediately, upgrading to
        128~GB when the price is favorable. Use a 2~TB primary NVMe and preserve
        the second slot.
  \item In early 2027, build the server only after deciding whether its first
        GPU is disposable. A 16~GB RTX 5060 Ti is a valid learning purchase near
        its reference price, but not a durable large-model solution.
  \item If budget permits, skip directly to a modern 16-core, 128~GB Linux
        server with a 24--32~GB NVIDIA GPU, two NVMe tiers, an upgrade-ready
        chassis/PSU, 2.5~GbE with a 10~GbE path, and a managed UPS.
  \item Prefer one larger-memory GPU for one larger model. Add a second GPU for
        separate services or measured sharding needs, not because aggregate
        VRAM appears additive on paper.
  \item Defer Threadripper and professional 48+~GB GPUs until NavKit campaigns,
        factor graphs, vision data, or multi-user services produce a documented
        requirement.
\end{enumerate}

This path delivers immediate development improvement, starts local inference
without pretending an entry card is permanent, and preserves a clean evolution
toward a balanced private engineering server.

\printbibliography
\end{document}
